\documentclass[preprint,12pt,authoryear]{elsarticle}

\usepackage{iftex}
\ifTUTeX                          % XeTeX/LuaTeX (tectonic): Unicode fonts, so ı İ ş ğ can be typed directly
  \usepackage{fontspec}           % Latin Modern by default
\else                             % pdfLaTeX
  \usepackage[T1]{fontenc}
  \usepackage[utf8]{inputenc}
  \usepackage{lmodern}
\fi
\usepackage{amsmath,amssymb}
\usepackage{booktabs}
\usepackage{array}
\usepackage{microtype}
\usepackage{graphicx}
\usepackage[hidelinks]{hyperref}

\graphicspath{{figures/}}

\journal{(journal to be decided)}

\begin{document}

\begin{frontmatter}

% Title to be written.
\title{TITLE}

\author[ytu]{H. O\u{g}uz Bolat\corref{cor}}
\ead{oguz.bolat@std.yildiz.edu.tr}
\cortext[cor]{Corresponding author}
\affiliation[ytu]{organization={Department of Geomatics Engineering, Y{\i}ld{\i}z Technical University},
  city={Istanbul}, country={T\"urkiye}}

\begin{abstract}
% Abstract.
\end{abstract}

\begin{keyword}
% Keywords, separated by \sep.
\end{keyword}

\end{frontmatter}

\section{Introduction}
% Motivation: sparse networks; why a generative model (neighbours explain
% only part of a station's series).

\section{Data}
% California cube; T\"urkiye cube; splits; held-out stations.

\section{Method}
% Problem; network thinning; conditional diffusion model.

\section{Experiments}
% A: held-out stations. B: planted fault-slip events. C: T\"urkiye and
% Kahramanmara\c{s}.

\section{Results}

\section{Discussion}

\section{Conclusions}

\section*{Data availability}
% GNSS time series: Nevada Geodetic Laboratory \citep{blewitt2018}.

\bibliographystyle{elsarticle-harv}
\bibliography{references}

\end{document}
